\chapter{Instalación de Arch Linux}

\newpage
\section{Conectar a Wi-Fi y Verificar Entorno}
El primer paso requiere establecer una conexión activa a internet mediante el entorno interactivo de instalación y comprobar la arquitectura de la placa madre.

\begin{enumerate}
\item[$\bullet$] \textbf{Verificación de Arquitectura EFI}: para corroborar que el sistema inició en modo UEFI, listamos las variables de firmware correspondientes.
\begin{lstlisting}
ls /sys/firmware/efi/efivars
\end{lstlisting}
\item[$\bullet$] \textbf{Conexión a la Red Inalámbrica}: utilizamos la herramienta interactiva \texttt{iwctl} para autenticar la tarjeta de red y conectarnos al punto de acceso:
\begin{lstlisting}
iwctl
# Dentro del entorno iwctl:
station wlan0 scan
station wlan0 get-networks
station wlan0 connect NombreRed
quit
\end{lstlisting}
\end{enumerate}

\begin{nota}[Instalación de Arch en Lenovo Ideapad 100]
\begin{enumerate}
\item[$\bullet$] \textbf{Solución de Errores Específicos de Hardware (Módulos Realtek)}
En laptops Lenovo IdeaPad 100 o equipos con tarjetas inalámbricas Realtek conflictivas (como la \texttt{rtl8723be}), es necesario forzar la antena correcta para evitar pérdidas de señal:
\begin{lstlisting}
modprobe -r rtl8723be
modprobe rtl8723be ant_sel=2
ip link set wlan0 up
\end{lstlisting}
\end{enumerate}
\end{nota}

\begin{enumerate}
\item[$\bullet$] \textbf{Verificación de Conectividad}
Comprobamos el envío y recepción de paquetes ICMP para asegurar el enlace:
\begin{lstlisting}
ping -c 3 google.com
\end{lstlisting}
\end{enumerate}

\newpage

\section{Particionado y Formateo del Disco}
Trabajaremos sobre el esquema del disco físico principal indexado como \texttt{/dev/sda}.

\begin{enumerate}
    \item [$\bullet$] \textbf{Limpieza de la Tabla de Particiones}
    Eliminamos cualquier estructura previa de almacenamiento y creamos una nueva tabla de particiones de tipo \textbf{GPT} (GUID Partition Table):
    \begin{lstlisting}
    parted /dev/sda
    # Dentro del entorno parted:
    mklabel gpt
    quit
    \end{lstlisting}
    \item [$\bullet$] \textbf{Estructura de Particiones en \texttt{cfdisk}}
    Accedemos a la herramienta visual \texttt{cfdisk} seleccionando la tabla de particiones de tipo \textbf{GPT}:
    \begin{lstlisting}
    cfdisk /dev/sda
    \end{lstlisting}

    Se define la siguiente distribución del almacenamiento:
    \begin{itemize}
        \item \textbf{\texttt{/dev/sda1}:} Nueva $\rightarrow$ $1\text{ GB}$ $\rightarrow$ Type: \textit{EFI System}.
        \item \textbf{\texttt{/dev/sda2}:} Nueva $\rightarrow$ $4\text{ GB}$ $\rightarrow$ Type: \textit{Linux swap}.
        \item \textbf{\texttt{/dev/sda3}:} Nueva $\rightarrow$ Resto del disco $\rightarrow$ Type: \textit{Linux filesystem}.
    \end{itemize}

    \item [$\bullet$] \textbf{Formateo de Particiones}
    Asignamos los sistemas de archivos correspondientes a cada bloque estructurado:
    \begin{lstlisting}
    # Particion EFI (FAT32)
    mkfs.fat -F 32 /dev/sda1

    # Particion Swap (Intercambio)
    mkswap /dev/sda2
    swapon /dev/sda2

    # Particion Raiz (ext4)
    mkfs.ext4 /dev/sda3
    \end{lstlisting}
\end{enumerate}

\begin{nota}[]
    Aplicar cambios escribiendo \texttt{write}, confirmar con \texttt{yes} y presionar \texttt{Quit} para salir.
\end{nota}

\newpage

\section{Instalación del Sistema Base}
Preparamos los puntos de montaje e inyectamos el núcleo del sistema operativo.

\begin{enumerate}
\item [$\bullet$] \textbf{Montaje Inicial de Unidades}
Montamos la raíz e inicializamos el directorio interno para la partición de arranque:
\begin{lstlisting}
mount /dev/sda3 /mnt
mkdir -p /mnt/boot
mount /dev/sda1 /mnt/boot
\end{lstlisting}

\item [$\bullet$] \textbf{Instalación mediante \texttt{pacstrap}}
Descargamos los paquetes esenciales en la raíz montada:
\begin{lstlisting}
pacstrap -K /mnt base linux linux-firmware nano
\end{lstlisting}

\item [$\bullet$] \textbf{Generación del Archivo de Puntos de Montaje (\texttt{fstab})}
Registramos los identificadores únicos de las particiones para los arranques subsecuentes:
\begin{lstlisting}
genfstab -U /mnt >> /mnt/etc/fstab
\end{lstlisting}
\end{enumerate}

\section{Configuración del Sistema Interno (\texttt{chroot})}

\begin{enumerate}
\item [$\bullet$] \textbf{Acceso al Entorno Virtualizado}
Ingresamos formalmente al disco duro del equipo para estructurar las directivas de usuario y localización.
\begin{lstlisting}
arch-chroot /mnt
\end{lstlisting}

\item [$\bullet$] \textbf{Zona Horaria y Reloj del Sistema}
Establecemos el huso horario correspondiente a Chile continental y sincronizamos los relojes de hardware:
\begin{lstlisting}
ln -sf /usr/share/zoneinfo/America/Santiago /etc/localtime
hwclock --systohc
\end{lstlisting}

\item [$\bullet$] \textbf{Localización e Idioma}
Editamos las plantillas generadoras de idioma del sistema:
\begin{lstlisting}
nano /etc/locale.gen
\end{lstlisting}

\item [$\bullet$] Descomentar desmarcando el caracter '\#' en la linea: 
\begin{lstlisting}
es_CL.UTF-8 UTF-8
\end{lstlisting}

\item [$\bullet$] \textbf{Compilación de Idiomas}
Compilamos los idiomas habilitados y fijamos la variable de entorno global:
\begin{lstlisting}
locale-gen
echo "LANG=es_CL.UTF-8" > /etc/locale.conf
\end{lstlisting}
\end{enumerate}

\section{Configuración del Sistema}

\begin{enumerate}
\item [$\bullet$] \textbf{Configuración del Teclado en Consola (TTY)}
Para que la distribución del teclado en modo texto reconozca la disposición en español:
\begin{lstlisting}
echo "KEYMAP=es" > /etc/vconsole.conf
\end{lstlisting}

\item [$\bullet$] \textbf{Identidad de Red (Hostname)}
Asignamos un nombre único a la máquina en la red local y asignamos la contraseña maestra de \texttt{root}:
\begin{lstlisting}
echo "Arch-Lenovo" > /etc/hostname
passwd
\end{lstlisting}

\item [$\bullet$] \textbf{Gestión de Redes Definitiva}
Instalamos los demonios correspondientes para asegurar la conectividad posterior:
\begin{lstlisting}
pacman -S networkmanager iw wpa_supplicant
systemctl enable NetworkManager
\end{lstlisting}

\item [$\bullet$] \textbf{Instalación y Despliegue del Cargador de Arranque (GRUB)}
Instalamos el gestor binario de arranque y escribimos sus bloques en la partición EFI activa:
\begin{footnotesize}
\begin{lstlisting}
pacman -S grub efibootmgr
grub-install --target=x86_64-efi --efi-directory=/boot --bootloader-id=GRUB
grub-mkconfig -o /boot/grub/grub.cfg
\end{lstlisting}
\end{footnotesize}


\item [$\bullet$] \textbf{Conclusión de la Fase Base y Reinicio}
Salimos del aislamiento del entorno virtual, desmontamos de forma recursiva y reiniciamos el hardware:
\begin{lstlisting}
exit
umount -R /mnt
reboot
\end{lstlisting}
\end{enumerate}

\begin{nota}[]
    Desconectar la unidad de almacenamiento USB (pendrive) en cuanto la pantalla cambie a negro.
\end{nota}

\newpage

\section{Configuración de Usuarios y Privilegios}
Tras el inicio autónomo del disco interno, accedemos con las credenciales de \texttt{root}.

\begin{enumerate}
\item [$\bullet$] \textbf{Conexión Inicial Mediante NetworkManager}
Invocamos la interfaz semi-gráfica para reactivar el enlace Wi-Fi sin dependencias de terceros:
\begin{lstlisting}
nmtui
# Seleccionar "Activar una conexion" -> Elegir Red
                                     -> Ingresar Password 
                                     -> Salir.
\end{lstlisting}

\item [$\bullet$] \textbf{Creación del Usuario Personal y Permisos \texttt{sudo}}
Instalamos la utilidad de elevación de privilegios y creamos el espacio de trabajo secundario:
\begin{lstlisting}
pacman -S sudo
useradd -m -G wheel -s /bin/bash guillermo
passwd guillermo
\end{lstlisting}

\item [$\bullet$] Para otorgar capacidades de administración al grupo del usuario, ejecutamos el editor seguro:
\begin{lstlisting}
EDITOR=nano visudo
# Buscar y descomentar removiendo el '#' de la linea:
# %wheel ALL=(ALL:ALL) ALL
\end{lstlisting}
Salimos de la sesión mediante \texttt{exit} e iniciamos sesión con el usuario definitivo (\texttt{guillermo}).
\end{enumerate}

\section{Interfaz Gráfica, Optimización de CPU y Entornos}

\begin{enumerate}
\item [$\bullet$] \textbf{Servidor Gráfico y Entorno de Escritorio Ligero (XFCE)}
Instalamos el servidor de despliegue Xorg junto al entorno tradicional optimizado y el gestor de accesos interactivo:
\begin{footnotesize}
\begin{lstlisting}
sudo pacman -S xorg-server xfce4 xfce4-goodies lightdm lightdm-gtk-greeter
sudo systemctl enable lightdm
\end{lstlisting}
\end{footnotesize}

\item [$\bullet$] \textbf{Optimización Avanzada de Energía para Procesadores Intel Celeron}
Para regular la frecuencia de los núcleos de forma dinámica, mitigar temperaturas elevadas y prolongar los ciclos de descarga en discos mecánicos (HDD):
\begin{lstlisting}
sudo pacman -S tlp tlp-rdw
sudo systemctl enable tlp.service
\end{lstlisting}

\item [$\bullet$] \textbf{Configuraciones del Entorno de Ventanas de Escritorio}
Para forzar la correcta traducción de caracteres y mapear las señales del teclado en el ecosistema X11 al español latinoamericano:
\begin{lstlisting}
localectl set-x11-keymap latam
\end{lstlisting}

\item [$\bullet$] \textbf{Despliegue del Navegador Web Principal}
Inyección del navegador nativo de código abierto:
\begin{lstlisting}
sudo pacman -S firefox
\end{lstlisting}
\end{enumerate}

\section{Alternativa Futura (Gestor de Ventanas Qtile)}
Para migrar a una maquetación en mosaico minimalista basada íntegramente en scripts dinámicos en lenguaje Python:

\begin{enumerate}
\item [$\bullet$] \textbf{Instalación de Componentes de Qtile y Terminal Acelerada}
\begin{lstlisting}
sudo pacman -S qtile python-psutil python-xdg python-iwlib alacritty
\end{lstlisting}

\item [$\bullet$] \textbf{Generación del Espacio de Trabajo en Python}
Estructuramos la jerarquía del archivo \texttt{config.py} heredando la configuración por defecto proporcionada por el sistema:
\begin{lstlisting}
mkdir -p ~/.config/qtile
cp /usr/share/doc/qtile/default_config.py ~/.config/qtile/config.py
\end{lstlisting}
\end{enumerate}

\begin{nota}{}
Tras cerrar sesión en LightDM, conmutar la sesión activa a ``Qtile'' e ingresar con las credenciales locales.
\end{nota}